\documentclass[dvipdfmx, uplatex]{jsarticle}
\usepackage[dvipdfmx]{graphicx}
\usepackage{amsmath}
\usepackage{amssymb}
\usepackage{amsfonts}
\usepackage{amsthm}
\usepackage{mathrsfs}
\usepackage{bm}
\usepackage{braket}
\usepackage{xr}
\externaldocument{exercise_lattice_random_walk}

\usepackage{silence}
\WarningFilter*{hyperref}{Token not allowed in a PDF string}
\usepackage[dvipdfmx]{hyperref}
\usepackage{pxjahyper}
\allowdisplaybreaks[4]

\usepackage[left=25truemm,top=20truemm,bottom=20truemm,right=25truemm]{geometry}

\newcommand{\bk}{\bm{k}}
\newcommand{\bx}{\bm{x}}

\theoremstyle{definition}
\newtheorem*{sol}{解答}

\title{解答：正方格子上のランダムウォークと拡散極限\\
\large ―― \texttt{exercise\_lattice\_random\_walk.tex} の解答 ――}
\author{}
\date{}

\begin{document}
\maketitle
\vspace{-6zh}

記号・式番号は問題ファイル \texttt{exercise\_lattice\_random\_walk.tex} のものを踏襲する。
\textbf{各問の解（採点用）}は太字（\boxed）で示す。

%======================================================================
\section{解答}
%======================================================================

\begin{sol}
\textbf{(1) 初期条件。}原点に局在ゆえ
\[
    \boxed{\,P_0(i,j)=\delta_{i,0}\,\delta_{j,0}\,}
\]
（時刻 $0$ で確率 $1$ が原点 $(0,0)$ に集中し，他の格子点では $0$ である状況）。

\par\smallskip
\textbf{(2) master 方程式。}時刻 $n+1$ に $(i,j)$ にいるのは，(a) 時刻 $n$ に $(i,j)$ にいて留まった
（確率 $1-p$），または (b) 時刻 $n$ に隣接 4 点のいずれかにいてそこから $(i,j)$ へ移ってきた
（各 $\tfrac p4$）場合だから
\[
    \boxed{\,P_{n+1}(i,j)=(1-p)\,P_n(i,j)+\frac{p}{4}\bigl[P_n(i+1,j)+P_n(i-1,j)+P_n(i,j+1)+P_n(i,j-1)\bigr]\,}.
\]
両辺を全格子点で和をとると右辺は $(1-p)\sum P_n+\tfrac p4\cdot4\sum P_n=\sum P_n$ ゆえ
$\sum_{i,j}P_{n+1}=\sum_{i,j}P_n$。$\sum_{i,j}P_0=1$ だから，すべての $n$ で $\boxed{\sum_{i,j}P_n(i,j)=1}$（確率保存）。

\par\smallskip
\textbf{(3) 特性関数。}(2) に $\sum_{i,j}e^{i\bk\cdot\bx}$ を掛けて和をとる。例えば
$\sum_{i,j}e^{i(k_xia+k_yja)}P_n(i+1,j)$ は $i'=i+1$ と置くと $e^{-ik_xa}\hat P_n(\bk)$ となる（他項も同様）から
\[
    \hat P_{n+1}(\bk)=\Bigl[(1-p)+\tfrac p4\bigl(e^{ik_xa}+e^{-ik_xa}+e^{ik_ya}+e^{-ik_ya}\bigr)\Bigr]\hat P_n(\bk),
\]
すなわち $\hat P_{n+1}=\lambda(\bk)\hat P_n$ で
\[
    \boxed{\,\lambda(\bk)=1-p+\frac{p}{2}\bigl(\cos k_xa+\cos k_ya\bigr)\,}.
\]
初期条件より $\hat P_0(\bk)=\sum_{i,j}\delta_{i,0}\delta_{j,0}e^{i\bk\cdot\bx}=1$ なので
\[
    \boxed{\,\hat P_n(\bk)=\lambda(\bk)^n\,}.
\]

\par\smallskip
\textbf{(4) 平均二乗変位。}1 ステップで $\Delta x=+a$（確率 $\tfrac p4$），$\Delta x=-a$（確率 $\tfrac p4$），
それ以外（$\pm y$ 移動または静止）は $\Delta x=0$ だから
$\langle\Delta x^2\rangle=\tfrac p4a^2+\tfrac p4a^2=\boxed{\dfrac{p a^2}{2}}$（$\langle\Delta y^2\rangle$ も同じ）。
各ステップは独立で平均変位 $0$ ゆえ分散は加法的で，$n$ ステップ後
\[
    \langle x^2\rangle_n=n\,\frac{pa^2}{2},\qquad
    \boxed{\,\langle x^2+y^2\rangle_n=n\,p\,a^2\,}.
\]
これは $n=t/\tau$ に比例する（拡散的な広がり）。
（特性関数からも $\langle x^2\rangle=-\partial_{k_x}^2\hat P_n|_{\bk=0}=-\partial_{k_x}^2\lambda^n|_0
=n\tfrac{pa^2}{2}$ と一致。）

\par\smallskip
\textbf{(5) 連続極限と diffusive scaling。}$\cos k_xa=1-\tfrac12(k_xa)^2+\cdots$ より
\[
    \lambda(\bk)=1-\frac{p a^2}{4}\,|\bk|^2+O\!\bigl((ka)^4\bigr).
\]
よって $n=t/\tau$ として
\[
    \hat P_n(\bk)=\lambda(\bk)^n=\exp\!\bigl[n\ln\lambda(\bk)\bigr]
    \;\xrightarrow{\ ka\ll1\ }\;\exp\!\Bigl[-\,\frac{t}{\tau}\,\frac{p a^2}{4}\,|\bk|^2\Bigr]
    =\exp\!\bigl[-D\,t\,|\bk|^2\bigr],
\]
ただし
\[
    \boxed{\,D=\frac{p\,a^2}{4\tau}\,}.
\]
$D$ を有限に保つには $a,\tau\to0$ で\textbf{$\dfrac{a^2}{\tau}$ を一定に保つ}必要がある
（$a\sim\sqrt{\tau}$，diffusive scaling）。$\tau\propto a$ や $\tau\propto a^3$ などでは $D\to\infty$ または $0$ となり
非自明な極限を与えない。

\par\smallskip
\textbf{(6) 連続極限の確率密度。}$e^{-Dt|\bk|^2}$ の逆 Fourier 変換（各次元 1 次元 Gauss 積分）より
\[
    P(x,y,t)=\int\frac{d^2k}{(2\pi)^2}\,e^{-i\bk\cdot\bx}\,e^{-Dt|\bk|^2}
    =\boxed{\,\frac{1}{4\pi D t}\,\exp\!\Bigl(-\frac{x^2+y^2}{4Dt}\Bigr)\,}
\]
（2 次元 Gauss 分布）。これは拡散方程式
\[
    \boxed{\,\frac{\partial P}{\partial t}=D\Bigl(\frac{\partial^2}{\partial x^2}+\frac{\partial^2}{\partial y^2}\Bigr)P
    =D\,\nabla^2 P\,}
\]
を満たす（$\hat P=e^{-Dt|\bk|^2}$ が $\partial_t\hat P=-D|\bk|^2\hat P$ を満たすことから）。
連続極限の平均二乗変位は Gauss 分布より $\langle x^2\rangle=\langle y^2\rangle=2Dt$ ゆえ
\[
    \boxed{\,\langle x^2+y^2\rangle=4Dt\,}.
\]
これは (4) の $\langle x^2+y^2\rangle_n=npa^2=\dfrac{t}{\tau}pa^2$ に $D=\dfrac{pa^2}{4\tau}$（(5)）を入れた
$4Dt$ と一致する。
\end{sol}

\end{document}
